% Version 0.26.  Complete finite-character rank-one Selmer suppliers.
\section{Finite character descent and rank-one vector vanishing}
\label{sec:rank-one-cohomology}

\status{RV1--RV7 close the literal suppliers \texttt{SoFin} and
\texttt{SoV} on every original algebraic closure and at every prime.
The proof uses the already reconstructed finite duality, local-K and
\texttt{LemmaK 31} suppliers.  It includes the finite-extension capture
of an algebraic-valued cocycle and the passage from a finite-length
torsion Selmer module to a vector coboundary.  RV1 also proves the
exact normalized-length transport needed when an auxiliary coefficient
field is used in the symplectic count.  The automorphy, rigidity,
counting and other cohomological packages in that count remain separate.}

Fix a prime $\ell$, a number field $k$, and its specified algebraic
closure $\Omega$.  Write $G_k=\operatorname{Gal}(\Omega/k)$ with its
profinite topology and $\chi=\chi_{k,\ell}:G_k\to\mathbb Z_\ell^\times$
for the actual cyclotomic character.  A coefficient ring $O$ will be
a discrete valuation ring, with its given $\mathbb Z_\ell$-algebra
structure, finite as a $\mathbb Z_\ell$-module.  As in the source,
a discrete valuation ring here is not a field.  Put
\[
 A_O=O\otimes_{\mathbb Z_\ell}(\mathbb Q_\ell/\mathbb Z_\ell),
 \qquad
 e_O=\operatorname{length}_O(O/\ell O).
 \tag{RV.1}
\]
Give $A_O$ the discrete topology.  For a character
$\psi:G_k\to O^\times$ with open kernel and an integer $e\ge1$, set
\[
 \alpha(g)=\psi(g)^{-1}\chi_O(g)^{-e},\qquad
 \rho(g)a=\alpha(g)a.
 \tag{RV.2}
\]
Here $\chi_O$ is obtained from the original $\chi$ by the specified
coefficient map; in particular the finite character is inverted.

For a valuation subring $A\subset\Omega$, let $D_A$ and $I_A$ denote
the actual decomposition and inertia subgroups of $G_k$.  Define
\[
 \begin{aligned}
 \mathscr D_\ell(k,\Omega)&=
 \{I_A:A\subset\Omega\text{ a valuation subring}\}
 \ \cup\
 \{D_A:A(\ell)<1\},\\
 \mathcal S(k,\Omega;\alpha)&=
 \bigcap_{U\in\mathscr D_\ell(k,\Omega)}
 \ker\!\left(H^1_{\rm cts}(G_k,A_O(\alpha))
                  \longrightarrow H^1_{\rm cts}(U,A_O(\alpha))\right).
 \end{aligned}
 \tag{RV.3}
\]
Thus a class in (RV.3) is represented by a continuous cocycle $z$
such that, for each indicated subgroup $U$, there is an $m_U\in A_O$
with $z(g)=(\alpha(g)-1)m_U$ for all $g\in U$.
These are precisely the source's
\texttt{LocSel rho top (shaFam k Omega ell)} and
\texttt{Chr.V ell e psi}; the latter is the extended
$O$-length of (RV.3).
No finite set of valuation subrings is substituted in this definition.

\subsection{Coefficients, normalized lengths and the finite duality join}

\begin{proposition}[RV1: exact coefficient extension and normalized length]
\label{claim:RV1}
The map $\mathbb Z_\ell\to O$ is injective, $O$ is finite free over
$\mathbb Z_\ell$, and $e_O=v_O(\ell)$ is a positive finite integer.
With $K=\operatorname{Frac}O$, there is a natural identification
\[
             A_O\simeq K/O,\qquad
 \operatorname{length}_O A_O[a]=v_O(a)\quad(0\ne a\in O).
 \tag{RV.4}
\]
Let $O\hookrightarrow O'$ be an injective compatible map of such
coefficient rings.  Put $b=v_{O'}(\varpi)$ for a uniformizer
$\varpi$ of $O$.  For every profinite group, every discrete
$O$-module representation, and every fixed family of local
subgroups, extension of coefficients gives the natural isomorphism
\[
 O'\otimes_O\mathcal S_{\mathscr D}(G,M)
       \simeq\mathcal S_{\mathscr D}(G,O'\otimes_O M).
 \tag{RV.5}
\]
For the rank-one coefficients in (RV.2), the right side uses the
transported character and the same original Galois and local groups.
Their extended lengths satisfy
\[
 \operatorname{length}_{O'}\mathcal S'
       =b\,\operatorname{length}_O\mathcal S,\qquad
 e_{O'}=b e_O.
 \tag{RV.6}
\]
Consequently finiteness is equivalent on the two sides, and, whenever
the lengths are finite,
\[
       \frac{\operatorname{length}_{O'}\mathcal S'}{e_{O'}}
          =\frac{\operatorname{length}_O\mathcal S}{e_O}.
 \tag{RV.7}
\]
The assertions apply to the source's level-$\ell N$ module
\texttt{Chr.V'} as well as to \texttt{Chr.V}.
\end{proposition}

\begin{proof}
A nonzero kernel of $\mathbb Z_\ell\to O$ would be its maximal
ideal, because $O$ is a domain.  Then $O$ would be a finite-dimensional
domain over $\mathbb F_\ell$, hence a field, contrary to the
coefficient hypothesis.  The map is therefore injective.
The finite torsion-free module $O$ over the principal ideal ring
$\mathbb Z_\ell$ is free.  The element $\ell$ is not a unit in $O$:
otherwise $O=\ell O$ and finite-module Nakayama over $\mathbb Z_\ell$
would give $O=0$.  Thus $0<v_O(\ell)<\infty$.
Factoring $\ell$ into a unit times a power of a uniformizer proves
the formula for $e_O$.

The finite-dimensional $\mathbb Q_\ell$-algebra $O[1/\ell]$ is a
domain and therefore a field: multiplication by any nonzero element
is injective and hence surjective.  It is $K$.
Tensor the exact sequence
$0\to\mathbb Z_\ell\to\mathbb Q_\ell\to
\mathbb Q_\ell/\mathbb Z_\ell\to0$ with the free module $O$.
This gives the first identification in (RV.4).
The $a$-torsion is $a^{-1}O/O$, isomorphic to $O/aO$.
In particular it has the asserted length, and multiplication by any
nonzero $a$ on $A_O$ is surjective.
The residue field of $O$ is finite, since it is a quotient of the
finite additive group $O/\ell O$.

The extension $O'$ is finite over $O$: any finite
$\mathbb Z_\ell$-generating set of $O'$ also generates it over $O$.
It is torsion-free, hence finite free.  The map is local.
Indeed a maximal ideal in an integral extension contracts to a
maximal ideal, so the contraction of the maximal ideal of $O'$ is
the unique maximal ideal of $O$.  Therefore $b$ is positive.

Choose an $O$-basis of $O'$.  Tensoring a discrete module with $O'$
then identifies it with a finite direct sum of copies of that module.
Continuous functions into this sum are exactly tuples of continuous
functions, so in every degree
\[
 O'\otimes_O C^q_{\rm cts}(G,M)
        \simeq C^q_{\rm cts}(G,O'\otimes_O M).
 \tag{RV.8}
\]
The natural map in (RV.8) commutes with the cochain differential
and every restriction map.  Finite freeness makes it exact.
In the chosen basis all restrictions act componentwise, so their
simultaneous kernels commute with this base change even when
$\mathscr D$ is infinite.  This proves (RV.5).
Tensor associativity identifies $O'\otimes_O A_O$ with $A_{O'}$;
the scalar action is the given transported action, not an arbitrary
isomorphic representation.

A finite composition series over $O$ has factors $O/\varpi O$.
After the exact base change, each factor is
$O'/\varpi O'$ and has $O'$-length $b$, proving the length formula
when the length is finite.
For completeness, infinite length is also preserved.  The element
$1\in O'$ is primitive in the finite free $O$-module $O'$:
if it belonged to $\varpi O'$, then $\varpi$ would be a unit in $O'$.
It can therefore be included in an $O$-basis.  It follows that
the map $M\to O'\otimes_O M$, $m\mapsto1\otimes m$, is split
injective over $O$.  If the latter module had finite $O'$-length,
it would have finite $O$-length, since the finite extension of
residue fields makes every simple $O'$-factor have finite $O$-length.
Then $M$ would have finite length as well.  This proves the extended
version of the first equality in (RV.6).
Finally write $\ell=u\varpi^{e_O}$ and apply $v_{O'}$ to obtain
the second equality.  Equation (RV.7) follows with both positive
denominators in place.  Replacing the local family by
$\mathscr D_{\ell N}$ in the same proof gives the level assertion.
\end{proof}

\begin{proposition}[RV2: the pure negative-twist strict Selmer module is finite]
\label{claim:RV2}
Suppose a number field $L\subset\Omega$ contains a primitive
$\ell$-th root of unity for odd $\ell$, or a primitive fourth
root for $\ell=2$.  For every $e\ge1$ and every original
$D:\texttt{LData}\ L\ \Omega\ \ell$,
\[
 \operatorname{length}_{\mathbb Z_\ell}
 \mathcal S(L,\Omega;\chi_{L,\ell}^{-e})
 \le
 \operatorname{length}_{\mathbb Z_\ell}
 H^2_{\rm cts}(G_{L,S_\ell},\mathbb Z_\ell(e+1))
 <\infty .
 \tag{RV.9}
\]
Here the coefficient on the left is $\mathbb Q_\ell/\mathbb Z_\ell$,
and the group on the right is the original compact-cochain group
defined using $D$.
In particular, if $n=[L:\mathbb Q]$ and
$\Delta_L=|\operatorname{disc}L|$, its left side is at most
\[
       \frac{\exp\!\big((2en)^{31}\log(2\Delta_L)\big)}
                       {\log\ell}.
 \tag{RV.10}
\]
\end{proposition}

\begin{proof}
The roots-of-unity hypothesis implies that $L$ is totally complex,
including when $\ell=2$.
The cyclotomic character is unramified away from $\ell$.
For instance, at a residue characteristic prime to $\ell$, the
polynomial $X^{\ell^r}-1$ has invertible derivative at every root,
so these roots lie in the maximal unramified local extension.
This also constructs the descended character in an \texttt{LData}
record; choose the finitely many primes above $\ell$ and one
extension of each valuation to $\Omega$.

On an inertia group away from $\ell$ the pure negative twist is
trivial.  A local coboundary there is therefore identically zero.
A cocycle with this property descends through the closed normal
subgroup generated by those inertia groups to $G_{L,S_\ell}$.
To see the normal-closure assertion directly, if a cocycle vanishes
on an inertia element acting trivially on the coefficient, its
cocycle identity makes it vanish on every conjugate; sums, inverses
and continuity then give the closed normal subgroup.
Conversely an inflated cocycle has all these inertia restrictions
zero.  Inflation is injective on $H^1$, by lifting the same
coboundary witness through the surjective quotient map.
At primes above $\ell$, its remaining condition is exactly the
strict local kernel defining $\operatorname{Sha}^1(D,e)$.
Changing a chosen valuation to a conjugate does not change vanishing
of a restriction: if $z|_U=\delta m$, then on $tUt^{-1}$ the
witness is $\alpha(t)m-z(t)$.
Thus the all-valuation module in (RV.9) is the original
$\operatorname{Sha}^1(D,e)$.

We recall the finite-to-integral duality argument explicitly.
Put $T_r=\mathbb Z/\ell^r(e+1)$ and
$A_r=(\mathbb Q_\ell/\mathbb Z_\ell)(-e)[\ell^r]$.
The Cartier dual of $T_r$ is $A_r$, since its twist is
$1-(e+1)=-e$.
LD3 gives perfect finite pairings
\[
 \operatorname{Sha}^2(L,T_r)\times
       \operatorname{Sha}^1(L,A_r)\longrightarrow\mathbb Q/\mathbb Z.
 \tag{RV.11}
\]
Reduction on $T_r$ is paired with inclusion on $A_r$.
This compatibility follows from the coefficient evaluation followed
by the same local invariant maps, so these are not unrelated choices
of perfect pairings.
KB6 transports the original decomposition groups, through their
chosen henselizations and completions, to the finite native local
fields of LK2--LK3.  It preserves roots of unity, cup products and
the actual restriction maps at every $r$, including at two.

KB1 supplies finite global $H^1$ at every coefficient level.
The local counterparts are finite by LK2.
Thus the relevant inverse systems of $H^1$ are Mittag--Leffler.
Continuous integral cochains are the inverse limit of the finite
cochains; each transition on cochains is surjective by lifting a
function on its finitely many clopen fibers.
The resulting $\varprojlim^1 H^1$ obstruction vanishes, globally
and locally.  Hence
\[
 \operatorname{Sha}^2(L,\mathbb Z_\ell(e+1))
       =\varprojlim_r\operatorname{Sha}^2(L,T_r).
 \tag{RV.12}
\]
For discrete torsion coefficients, continuous cochains commute with
the filtered union: their domains are compact and their images
finite.  There are finitely many local places in this presentation,
and filtered colimits are exact.  Therefore
\[
 \operatorname{Sha}^1(L,(\mathbb Q_\ell/\mathbb Z_\ell)(-e))
       =\varinjlim_r\operatorname{Sha}^1(L,A_r).
 \tag{RV.13}
\]
The compatible finite pairings identify the compact dual of
(RV.13) with (RV.12), a submodule of the actual integral $H^2$.
This is also the finite duality argument underlying (TW.12);
no Iwasawa or Ferrero--Washington theorem is used here.
KB8 makes that $H^2$ finite and bounds the logarithm of its
cardinality by the numerator of (RV.10).
A finite $\mathbb Z_\ell$-module of length $h$ has cardinality
$\ell^h$.  Duality preserves that length, proving (RV.9)--(RV.10).
\end{proof}

\subsection{The actual finite-character supplier}

\begin{theorem}[RV3: finite-character descent and literal \texttt{SoFin}]
\label{claim:RV3}
For every original $k,\Omega,\ell,O,\psi$ as in (RV.1)--(RV.2)
and every $e\ge1$,
\[
                    \texttt{Chr.V}\ \ell\ e\ \psi<\infty .
 \tag{RV.14}
\]
No assumption that the order of $\psi$ is prime to $\ell$ is made.
In particular the source's literal \texttt{SoFin Omega ell} holds
for every algebraic closure $\Omega$ of $\mathbb Q$ and every prime.
\end{theorem}

\begin{proof}
Let $N=\ell$ for odd $\ell$ and $N=4$ at two.  Inside the original
$\Omega$ take the fixed field
\[
 L=\Omega^{\,\ker\psi\cap\ker(\chi\bmod N)}.
 \tag{RV.15}
\]
Both kernels are open, so $L/k$ is finite Galois.
Its Galois group embeds in the product of two abelian character
images, so it is abelian.  The first kernel makes $\psi|_{G_L}=1$.
The second fixes every $N$-th root of unity in $\Omega$; hence $L$
contains a primitive such root.  It is totally complex and satisfies
the exact root hypothesis of RV2.
This is an extension of the original $k$ inside the original
$\Omega$, with the original embeddings.

Set $H=G_L$, $\Gamma=G_k$, $\Delta=\Gamma/H$, and $d=|\Delta|$.
There is an $h\in H$ with $\chi(h)^e\ne1$.
Indeed the mod-$\ell^r$ cyclotomic image of $G_L$ has order
$[L(\mu_{\ell^r}):L]\ge
\varphi(\ell^r)/[L:\mathbb Q]$, which tends to infinity.
The degree formula used here has a short direct proof.
The primitive $\ell^r$-th roots are the roots of
$P_r(X)=\sum_{j=0}^{\ell-1}X^{j\ell^{r-1}}$.
Modulo $\ell$, the polynomial $P_r(1+T)$ is
$T^{\ell^{r-1}(\ell-1)}$, and its constant coefficient is
$\ell$, not divisible by $\ell^2$.
Eisenstein's criterion makes it irreducible over $\mathbb Q$,
of degree $\varphi(\ell^r)$.
The image of $\chi|_H$ is therefore infinite, whereas the equation
$u^e=1$ has at most $e$ roots in $\mathbb Q_\ell$.
Put
\[
                       a=\chi(h)^e-1\in\mathbb Z_\ell\setminus\{0\}.
 \tag{RV.16}
\]
On $H$ the coefficient action is the pure negative twist.
Consequently
$B=A_O^H\subset A_O[a]$: the scalar
$\chi(h)^{-e}-1$ differs from $a$ by a unit.
RV1 gives
\[
             \operatorname{length}_O B
                     \le e_O\,v_\ell(a)<\infty.
 \tag{RV.17}
\]
This supplies the needed $H^0$ control before using any restriction
kernel bound.

We give that bound without averaging over $\Delta$.
If a cocycle $z$ restricts to a coboundary on $H$, subtract that
coboundary so that $z|_H=0$.
The identity $hg=g(g^{-1}hg)$ shows
$\alpha(h)z(g)=z(g)$ for $h\in H$.
Also $z(gh)=z(g)$, so $z$ descends to a cocycle of $\Delta$
with values in $B$.
Two such normalized cocycles differ by a coboundary precisely when
their witnesses differ by an element of $B$.
This proves the exact kernel identification
\[
 \ker\!\left(H^1(\Gamma,A_O(\alpha))
                  \longrightarrow H^1(H,A_O(\alpha))\right)
                      \simeq H^1(\Delta,B).
 \tag{RV.18}
\]
A normalized finite-group cocycle is determined by its values on
the $d-1$ nonidentity elements.  Therefore
$\operatorname{length}_O H^1(\Delta,B)
\le(d-1)\operatorname{length}_O B$.
This remains valid when $\ell$ divides $d$.

Restriction sends the strict module (RV.3) into the strict module
over $L$.  For the same valuation subring $A\subset\Omega$, the
local groups over $L$ are $H\cap D_A$ and $H\cap I_A$; the
condition $A(\ell)<1$ is unchanged.  Thus every original local
coboundary witness restricts to the required witness.
Since $\psi|_H=1$, RV1 applied to the pure twist gives
\[
 \operatorname{length}_O\mathcal S(L,\Omega;\alpha|_H)
       =e_O\,
          \operatorname{length}_{\mathbb Z_\ell}
                  \mathcal S(L,\Omega;\chi_{L,\ell}^{-e}).
 \tag{RV.19}
\]
Its right side is finite by RV2.
The image and kernel of the restriction of (RV.3) therefore both
have finite length.  More explicitly, (RV.17)--(RV.19) yield
\[
 \frac{\operatorname{length}_O\mathcal S(k,\Omega;\alpha)}{e_O}
 \le
 \operatorname{length}_{\mathbb Z_\ell}
                  \mathcal S(L,\Omega;\chi_{L,\ell}^{-e})
                    +(d-1)v_\ell(a).
 \tag{RV.20}
\]
The numerator on the left is proved finite before this real-valued
normalization is used.
This proves (RV.14) for an arbitrary original coefficient ring.
Setting $k=\mathbb Q$ and $O=c.O$ proves the exact \texttt{SoFin}
predicate.  The estimate (RV.20) is a sufficient alternative to
the source's sharper generator-count restriction estimate; no
unstated prime-to-$\ell$ restriction or new finiteness hypothesis
has been introduced.
\end{proof}

\subsection{Capturing the coefficient field of a continuous cocycle}

\begin{proposition}[RV4: finite-extension capture from Baire and the cocycle equation]
\label{claim:RV4}
Let $\Gamma$ be a compact Hausdorff topological group,
$K_0\subset\overline{\mathbb Q}_\ell$ a finite extension of
$\mathbb Q_\ell$, and let
$\theta,c:\Gamma\to\overline{\mathbb Q}_\ell$ satisfy
\[
 \theta(g)\in K_0,\qquad \theta(1)=1,\qquad
 c\text{ continuous},\qquad
 c(gh)=c(g)+\theta(g)c(h).
 \tag{RV.21}
\]
Then a finite extension $E\subset\overline{\mathbb Q}_\ell$
contains $K_0$ and every value of $c$.
Here $\overline{\mathbb Q}_\ell$ is the algebraic closure with
its extended valuation, not its completion.
\end{proposition}

\begin{proof}
We first construct a countable family of finite subfields
containing every finite subset of $\overline{\mathbb Q}_\ell$.
Choose an embedded algebraic closure $\overline{\mathbb Q}$ of
$\mathbb Q$ in $\overline{\mathbb Q}_\ell$.
It is countable and dense.
For the density assertion, let $\alpha$ be algebraic over
$\mathbb Q_\ell$ with monic minimal polynomial $f$ of degree $n$.
Approximate the coefficients of $f$ by rational numbers, retaining
the leading coefficient $1$, so that the resulting
$g\in\mathbb Q[X]$ satisfies $|g(\alpha)|<\varepsilon^n$.
This is possible because $\mathbb Q$ is dense in $\mathbb Q_\ell$
and there are finitely many coefficients.
All roots $\beta_1,\ldots,\beta_n$ of $g$ lie in the embedded
$\overline{\mathbb Q}$, and
$\prod_i|\alpha-\beta_i|=|g(\alpha)|$.
At least one $\beta_i$ is within $\varepsilon$ of $\alpha$.

For $\alpha\notin\mathbb Q_\ell$, choose such a $\beta$ satisfying
\[
 |\alpha-\beta|<
       \min_{\alpha'\ne\alpha}
                         |\alpha-\alpha'|,
 \tag{RV.22}
\]
where $\alpha'$ ranges over its $\mathbb Q_\ell$-conjugates.
Every automorphism of $\overline{\mathbb Q}_\ell$ fixing
$\mathbb Q_\ell$ is an isometry, by uniqueness of the extended
absolute value.  If it also fixes $\beta$, then
$|\sigma\alpha-\alpha|\le|\alpha-\beta|$, so (RV.22) forces
$\sigma\alpha=\alpha$.
Finite Galois theory gives $\alpha\in\mathbb Q_\ell(\beta)$.
This proves the needed Krasner assertion directly.
If $\alpha\in\mathbb Q_\ell$, no adjoining is needed.
The fields generated over $\mathbb Q_\ell$ by finite subsets of
the countable $\overline{\mathbb Q}$ consequently form a
countable family $(K_i)$, cofinal for finite subsets of
$\overline{\mathbb Q}_\ell$.
Each $K_0K_i$ is finite over $\mathbb Q_\ell$, complete for
its extended absolute value, and hence closed in
$\overline{\mathbb Q}_\ell$.

The cocycle identity gives $c(1)=0$ and
$c(g^{-1})=-\theta(g^{-1})c(g)$.
Thus
\[
                 H_i=\{g\in\Gamma:c(g)\in K_0K_i\}
 \tag{RV.23}
\]
is a subgroup: products and inverses stay in the indicated field.
It is closed by continuity of $c$.
The $H_i$ cover $\Gamma$.
The Baire property of a compact Hausdorff space now implies that
some $H_i$ has nonempty interior.
One can see the needed Baire assertion directly: if a countable
closed cover had all members with empty interior, successive
nonempty open sets could be chosen with compact closures nested
inside the previous open set and avoiding the next member.
Compactness would give a point in all those closures, outside
the proposed cover.

A subgroup with nonempty interior is open, and an open subgroup
of a compact group has finite index.
Choose finitely many representatives $g_1,\ldots,g_s$ for its
left cosets and let
\[
                    E=K_0K_i(c(g_1),\ldots,c(g_s)).
 \tag{RV.24}
\]
This is a finite extension because all its generators are algebraic.
For $g=g_jh$ with $h\in H_i$, the identity (RV.21) gives
$c(g)=c(g_j)+\theta(g_j)c(h)\in E$.
This proves the assertion.
The subgroup structure in (RV.23) is essential: the argument
does not assert that every compact subset of
$\overline{\mathbb Q}_\ell$ is contained in a finite extension.
\end{proof}

\subsection{From torsion Selmer finiteness to an exact vector coboundary}

\begin{proposition}[RV5: a finite-length quotient controls vector cocycles]
\label{claim:RV5}
Let $E/\mathbb Q_\ell$ be finite, $O_E$ its valuation ring,
$\alpha:\Gamma\to O_E^\times$ a continuous scalar action, and
$\mathscr D$ any family of subgroups of a topological group $\Gamma$.
Suppose
$\mathcal S_{\mathscr D}(\Gamma,(E/O_E)(\alpha))$
has finite $O_E$-length.
If a continuous $c:\Gamma\to E$ is an $\alpha$-cocycle and is
a coboundary on every $U\in\mathscr D$, with witnesses in $E$,
then there is an $m\in E$ with
\[
                         c(g)=(\alpha(g)-1)m
                         \quad\text{for every }g.
 \tag{RV.25}
\]
\end{proposition}

\begin{proof}
The quotient map $\pi:E\to E/O_E$ is continuous for the discrete
quotient topology, since its kernel $O_E$ is open.
A finite-length module of length $h$ is killed by $\ell^h$:
each simple composition factor is killed by the maximal ideal,
which contains $\ell$, and induction along the series gives
the assertion.  Choose $N$ such that $\ell^N$ kills the strict
Selmer module in the hypothesis.

For each $n\ge0$, reduce
$\ell^{-(n+N)}c$ modulo $O_E$.
It is a continuous cocycle in that strict module; its local
witness on $U$ is the reduction of $\ell^{-(n+N)}m_U$.
Multiplication by $\ell^N$ kills its class, so for some $y_n\in E$
\[
            \ell^{-n}c(g)-(\alpha(g)-1)y_n\in O_E
                         \quad\text{for all }g.
 \tag{RV.26}
\]
Put $m_n=\ell^n y_n$ and $r=|\ell|=\ell^{-1}<1$.
Then
\[
                 |c(g)-(\alpha(g)-1)m_n|\le r^n
                         \quad\text{for all }g.
 \tag{RV.27}
\]

If $\alpha(g)=1$ for every $g$, (RV.27) directly gives $c=0$,
and $m=0$ works.
Otherwise choose $g_0$ with $d_0=\alpha(g_0)-1\ne0$ and set
$m=c(g_0)/d_0$.
For $d_g=\alpha(g)-1$, the exact identity
\[
 c(g)-d_gm
   =\big(c(g)-d_gm_n\big)
        -\frac{d_g}{d_0}\big(c(g_0)-d_0m_n\big)
 \tag{RV.28}
\]
bounds its absolute value by
$r^n\max(1,|d_g/d_0|)$ for every $n$.
Letting $n$ tend to infinity proves (RV.25).
The sequence of witnesses $m_n$ need not be asserted to converge:
the single explicit witness $m=c(g_0)/d_0$ gives the conclusion.
\end{proof}

\begin{proposition}[RV6: the vector theorem on an actual finite coefficient field]
\label{claim:RV6}
Let $E\subset\overline{\mathbb Q}_\ell$ be finite.
Let $\psi:G_k\to\overline{\mathbb Q}_\ell^\times$ have open
kernel and all its values in $E$, and let $e\ge1$.
Suppose $\theta,c:G_k\to\overline{\mathbb Q}_\ell$ satisfy
\[
 \theta(g)\psi(g)\chi(g)^e=1,\qquad
 c\text{ continuous},\qquad c(g)\in E,\qquad
 c(gh)=c(g)+\theta(g)c(h).
 \tag{RV.29}
\]
Suppose also that $c$ is a coboundary on every original inertia
group, and on every original decomposition group above $\ell$,
with witnesses initially allowed in $\overline{\mathbb Q}_\ell$.
Then a witness $m$ in this same $E$ makes $c$ a global coboundary.
This is the complete finite-field statement \texttt{soVec}.
\end{proposition}

\begin{proof}
We record the coefficient-lattice facts needed here.
The finite extension $E$ has a unique extended absolute value,
is complete, and its valuation ring
$O_E=\{x\in E:|x|\le1\}$ is the integral closure of
$\mathbb Z_\ell$.
The last equality can also be checked from the minimal polynomial.
An integral element cannot have absolute value greater than one,
since then the leading term of a monic integral relation would
strictly dominate the others.
Conversely, if $|x|\le1$, all its conjugates have absolute value
at most one; the elementary symmetric functions show that its
monic minimal polynomial has coefficients in $\mathbb Z_\ell$.

For clarity, $O_E$ is finite free over $\mathbb Z_\ell$.
Choose a $\mathbb Q_\ell$-basis consisting of integral elements.
The traces of their products with any element of $O_E$ belong
to $\mathbb Z_\ell$.
Thus $O_E$ lies in the trace-dual lattice of that basis.
Nondegeneracy of the trace pairing in characteristic zero and
Noetherianity of $\mathbb Z_\ell$ make it a finite torsion-free,
hence free, module.
Its value group is discrete: the norm formula gives
$|x|^{[E:\mathbb Q_\ell]}
=|N_{E/\mathbb Q_\ell}(x)|\in\ell^{\mathbb Z}$.
The smallest positive additive valuation yields a uniformizer;
factoring by it shows that $O_E$ is a discrete valuation ring.
It is not a field because $0<|\ell|<1$.
It therefore has exactly the coefficient properties required in RV3,
and (RV.4) identifies $A_{O_E}$ with $E/O_E$.

Compactness of $G_k$ and openness of $\ker\psi$ make its quotient
finite.  Every value of $\psi$ is a root of unity and hence has
absolute value one.  Thus $\psi$ is a character
$\psi_E:G_k\to O_E^\times$ with exactly the same kernel.
The identity in (RV.29) uniquely determines
\[
                    \theta=\psi_E^{-1}\chi_{O_E}^{-e};
 \tag{RV.30}
\]
in particular $\theta$ takes values in $O_E^\times$ and is continuous.
The map $c$, with values in the subspace $E$, is continuous as
an $E$-valued map.
RV3 supplies finite length for the exact strict module
$\mathcal S(k,\Omega;(E/O_E)(\theta))$.

Each local witness can be taken in $E$ without enlarging it.
For a local subgroup $U$, if $\theta|_U=1$, its coboundary
condition already says $c|_U=0$, and the witness $0$ works.
Otherwise choose $g_0\in U$ with $\theta(g_0)\ne1$.
Any witness must equal
\[
                         \frac{c(g_0)}{\theta(g_0)-1}\in E.
 \tag{RV.31}
\]
It therefore gives the entire local coboundary identity in $E$.
This argument treats each original subgroup directly; there is no
attempt to adjoin infinitely many independently selected witnesses.
Apply RV5 with the unchanged family
$\mathscr D_\ell(k,\Omega)$ and the coefficient identification
above.  The resulting $m\in E$ proves the assertion.
\end{proof}

\begin{theorem}[RV7: the literal rank-one suppliers on every original closure]
\label{claim:RV7}
For every original algebraic closure $\Omega$ of $\mathbb Q$
and every prime $\ell$,
\[
       \texttt{SoV}\ \Omega\ \ell
              \quad\text{and}\quad
       \texttt{SoFin}\ \Omega\ \ell
 \tag{RV.32}
\]
hold.  The vector statement in fact holds over every number
field $k$ in the full generality of \texttt{soBar}.
Under an injective extension of coefficient rings, the original
\texttt{Vlen} is finite if and only if its transported version is
finite, and its normalized \texttt{Vn} is unchanged.
More generally, the same conclusion holds for two coefficient rings
inside the prescribed $\overline{\mathbb Q}_\ell$ whose Dirichlet
characters have equal values there.  This coefficient assertion
allows every $e\ge0$ and is the literal \texttt{coefIndep} statement.
\end{theorem}

\begin{proof}
RV3 gives \texttt{SoFin}.
For the vector statement, take exactly the original
$k,\Omega,\ell,\psi,e,\theta,c$ and its original local witnesses
in the definition of \texttt{soBar}.
The finite image of $\psi$ generates a finite
$K_0\subset\overline{\mathbb Q}_\ell$.
Since $\theta(g)=(\psi(g)\chi(g)^e)^{-1}$, every value of
$\theta$ lies in $K_0$ and $\theta(1)=1$.
RV4 therefore supplies a finite $E\supset K_0$ containing all
values of $c$.
RV6 applies on this same $E$ and produces a global witness.
Regard that witness as an element of the originally prescribed
$\texttt{PadicAlgCl}\ \ell$; this is the exact conclusion of
\texttt{soBar}.  Specializing $k=\mathbb Q$ is precisely
\texttt{SoV Omega ell}.

For the last statement, the Dirichlet character transported by an
injective coefficient map gives the transported Galois character
by composition with the same cyclotomic restriction map.
The source's \texttt{Vlen} is the length of the level-$\ell N$
module \texttt{Chr.V'}, and \texttt{Vn} is that length divided by
$e_O$.  Apply RV1 to that unchanged local family.
Both lengths are finite together, and (RV.7) is the required
exact normalized equality in the finite case.  If both lengths
are infinite, the source's convention $\infty.\texttt{toNat}=0$
makes both displayed real normalizations zero; this convention
does not assert a finite cohomological bound.
For two possibly nonnested coefficient rings, take the compositum
of their fraction fields inside $\overline{\mathbb Q}_\ell$.
It is finite over $\mathbb Q_\ell$ by RV1.
Its valuation ring, constructed in the proof of RV6, contains both
original rings because their elements are integral over
$\mathbb Z_\ell$.  The two transported characters there are
identical by the given equality of their original values.
Apply RV1 to the two inclusions and compare through this common
ring.  This proves the full \texttt{coefIndep} assertion, without
changing the original $\Omega$, level or local family.
In particular this step returns a
bound from an auxiliary coefficient field to the original ring;
it does not merely produce an estimate on the auxiliary ring.
\end{proof}

\paragraph{Classical input and scope.}
The new local-field background used above is the uniqueness of the
extended absolute value on a finite extension of $\mathbb Q_\ell$,
its completeness and the norm formula.  See J.\,S.~Milne,
\emph{Algebraic Number Theory}, version 3.08 (2020),
Theorem 7.38 and Corollary 7.40, pp.~123--124,
\url{https://www.jmilne.org/math/CourseNotes/ANT.pdf}.
The comparison in (RV.22) is proved here; it is also Proposition
7.60 of those notes, pp.~131--132.
The elementary prime-power cyclotomic degree formula is proved
by Eisenstein's criterion in RV3; compare Proposition 6.2(a)
of those notes, p.~96.
The finite Galois correspondence, elementary cochain algebra,
finite-module structure over a discrete valuation ring, and the
Baire argument were used in their explicit forms above.
LD3, LK2--LK3 and KB1/KB6/KB8 retain their own previously recorded
classical theorem leaves.

No Iwasawa control theorem, arbitrary compact-set capture principle,
vanishing of an unconstructed automorphic representation, or
prime-to-$\ell$ hypothesis on the finite character is an input.
The two predicates in (RV.32) are now actual supplied leaves in
the symplectic count.  Its \texttt{CorX}, other cohomology package,
finite algebra models, automorphy, rigidity and trace-counting
requirements are separate.  This chapter alone does not assert
\texttt{CountII}, \texttt{CountIII}, \texttt{PointHV}, RH or GRH.
